% =============================================================================
%  Brick Routing Algorithm -- mathematical description.
%  Include-ready LaTeX section: no documentclass, no preamble, no bibliography.
% =============================================================================

\section{Brick Routing Algorithm}
\label{sec:brick-routing-algorithm}

This section describes the Brick routing algorithm as a mathematical decision
rule. The router receives a user query, estimates which capabilities the query
requires, estimates its difficulty, and selects one model from a fixed pool by
minimizing a cost-aware capability-distance objective. The implementation used
in the experiments is deterministic at routing time: stochasticity appears only
in the offline hyperparameter search used to select the locked production
configuration.

\subsection{Model Pool and Notation}
\label{sec:brick-notation}

Let
\[
  \mathcal{M}=\{m_1,\ldots,m_K\}
\]
be the candidate model pool. In the Dataset~A experiments,
\[
  \mathcal{M} =
  \{\texttt{qwen3.5-9b},\texttt{deepseek-v4-flash},\texttt{kimi2.6}\}.
\]
Let
\[
  \mathcal{C}=\{c_1,\ldots,c_D\}
\]
be the ordered capability basis. The current six-dimensional basis is
\[
\begin{split}
\mathcal{C}=\{&
\texttt{coding},\texttt{creative\_synthesis},
\texttt{instruction\_following},\\
&\texttt{math\_reasoning},\texttt{planning\_agentic},
\texttt{world\_knowledge}\}.
\end{split}
\]
Each model \(m\) has a non-negative scalar cost \(a_m\). In the reported
experiments these normalized costs are
\[
  a_{\texttt{qwen}}=0.10,\qquad
  a_{\texttt{ds4}}=0.40,\qquad
  a_{\texttt{kimi}}=0.60.
\]

For each model and capability, the router stores a skill value
\[
  s_{m,c}\in(0,1),
\]
interpreted as the empirical probability that model \(m\) succeeds on work
requiring capability \(c\). These values form a skill matrix
\[
  S\in(0,1)^{K\times D}.
\]

\subsection{Query Representation}
\label{sec:brick-query-representation}

For a query \(x\), a capability classifier returns a probability distribution
over the capability basis:
\[
  p(x) = (p_1,\ldots,p_D), \qquad
  p_c \ge 0,\qquad \sum_{c=1}^{D} p_c = 1.
\]
The probability \(p_c\) is not a hard label. It is a soft allocation of the
query's required work across capabilities.

A complexity classifier returns a label
\[
  y\in\{\mathrm{easy},\mathrm{medium},\mathrm{hard}\}
\]
and confidence \(\rho\in[0,1]\). Labels are mapped to base difficulty levels
\[
  \tau_{\mathrm{easy}}=0.55,\qquad
  \tau_{\mathrm{medium}}=0.72,\qquad
  \tau_{\mathrm{hard}}=0.88.
\]
The effective query difficulty is confidence-blended toward the medium prior:
\[
  \tau_q = \rho\,\tau_y + (1-\rho)\,\tau_{\mathrm{medium}}.
\]
If the label is missing or invalid, the router sets \(\rho=0\), so that
\(\tau_q=\tau_{\mathrm{medium}}\).

The router works in log-odds space. Define
\[
  \operatorname{logit}(u)=\log\frac{u}{1-u},
\]
with probabilities clipped to the interval \([\epsilon_{\min},\epsilon_{\max}]\)
before applying the logit. The production clipping defaults are
\[
  \epsilon_{\min}=0.02,\qquad \epsilon_{\max}=0.98.
\]

\subsection{Skill-Vector Calibration}
\label{sec:brick-skill-calibration}

In calibrated experiments, the skill matrix is estimated from a development
set with known per-model correctness labels. Let \(z_{i,m}\in\{0,1\}\) denote
whether model \(m\) correctly solves training query \(i\), and let \(p_{i,c}\)
be the capability probability assigned to capability \(c\) for that query.
The empirical skill estimate is the probability-weighted success rate
\[
  \widehat{s}_{m,c}
  =
  \frac{\sum_{i\in\mathcal{D}_{\mathrm{cal}}} p_{i,c} z_{i,m}}
       {\max\left(\sum_{i\in\mathcal{D}_{\mathrm{cal}}} p_{i,c},\delta\right)},
\]
where \(\delta\) is a small numerical constant. The estimate is clipped to
\([0.02,0.98]\) before routing:
\[
  s_{m,c}=\operatorname{clip}(\widehat{s}_{m,c},0.02,0.98).
\]
This is not model fine-tuning. No target model weights are updated. Only the
router's per-model capability table is calibrated from observed success labels.

\subsection{Base Routing Objective}
\label{sec:brick-base-objective}

The scalar difficulty \(\tau_q\) is transformed into an effective complexity
logit
\[
  z_q = b + \mu \operatorname{logit}(\tau_q),
\]
where \(b\) is the complexity bias and \(\mu\) is the complexity multiplier.
The query requirement vector is
\[
  r_{q,c}=p_c z_q.
\]

For each model, the comparable model vector is
\[
  v_{m,c}=p_c \operatorname{logit}(s_{m,c}).
\]
The multiplication by \(p_c\) has two effects. First, it focuses the distance
on the capabilities relevant to the current query. Second, it keeps unrelated
capabilities from dominating the score simply because a model is strong in an
irrelevant dimension.

The router distinguishes under-capacity from over-capacity:
\[
  u_{m,c} = \max(0,r_{q,c}-v_{m,c}),
  \qquad
  o_{m,c} = \max(0,v_{m,c}-r_{q,c}).
\]
Under-capacity means the model appears weaker than the query requirement.
Over-capacity means the model appears stronger than required; this is not a
correctness problem, but it can be an economic problem when the stronger model
is expensive.

The capability distance for model \(m\) is
\[
  D_m
  =
  \sqrt{
    \sum_{c=1}^{D}
    \left(
      u_{m,c}^{2}+\lambda o_{m,c}^{2}
    \right)
  },
\]
where \(\lambda\ge0\) controls the penalty for over-routing. The final score is
\[
  J_m = D_m + \beta a_m,
\]
where \(\beta\ge0\) controls the cost penalty. The selected model is
\[
  m^\star = \arg\min_{m\in\mathcal{M}} J_m.
\]

For debugging and deterministic tie-breaking, the router also computes
\[
  E_m=\sum_{c=1}^{D} p_c s_{m,c},
\]
the expected success under the query capability mixture. If two models have
scores within the tie tolerance \(\varepsilon_{\mathrm{tie}}\), the router
chooses the model with higher \(E_m\); if this is also tied, it chooses the
lower-cost model.

\subsection{User Preference Knob}
\label{sec:brick-knob}

The router exposes a continuous preference knob
\[
  r\in[-1,1].
\]
The three named profiles are
\[
  \mathrm{min}\equiv -1,\qquad
  \mathrm{neutral}\equiv 0,\qquad
  \mathrm{max}\equiv +1.
\]
The intent is asymmetric:
\begin{itemize}
  \item \textbf{min}: minimize cost, accepting lower capability;
  \item \textbf{neutral}: use the calibrated quality--price compromise;
  \item \textbf{max}: maximize answer quality, strongly relaxing cost pressure.
\end{itemize}

Let \(\alpha>0\) be the preference power. Define
\[
  r_+ = \max(r,0)^\alpha,
  \qquad
  r_- = \max(-r,0)^\alpha.
\]
The effective routing parameters are
\begin{align}
  \mu(r)
  &=
  \mu_0
  \exp\left(
    r_+\log A_{\mu}^{+}
    + r_-\log A_{\mu}^{-}
  \right),\\
  b(r)
  &=
  b_0 + r_+ A_b^{+} + r_- A_b^{-},\\
  \beta(r)
  &=
  \beta_0
  \exp\left(
    -r_+\log A_{\beta}^{+}
    + r_-\log A_{\beta}^{-}
  \right),\\
  \lambda(r)
  &=
  \lambda_0
  \exp\left(
    -r_+\log A_{\lambda}^{+}
    + r_-\log A_{\lambda}^{-}
  \right).
\end{align}
Here \((\mu_0,b_0,\beta_0,\lambda_0)\) are the neutral parameters.
The \(+\) constants control the max-performance direction; the \(-\)
constants control the economy direction.

At \(r=0\), \(r_+=r_-=0\), therefore
\[
  \mu(0)=\mu_0,\quad b(0)=b_0,\quad
  \beta(0)=\beta_0,\quad \lambda(0)=\lambda_0.
\]
At \(r=1\), the router increases the query requirement and sharply reduces
cost and over-capacity penalties. This makes the score favor stronger models
even when they are more expensive. At \(r=-1\), the router compresses the
requirement and increases the economic penalties, making cheaper models win
unless the distance penalty is overwhelming.

\begin{table}[ht]
\centering
\caption{Locked production math configuration for the aggressive preference
knob.}
\label{tab:brick-knob-defaults}
\small
\begin{tabular}{l r}
\toprule
\textbf{Parameter} & \textbf{Value} \\
\midrule
\(\mu_0\) (\texttt{complexity\_mu}) & 1.07 \\
\(b_0\) (\texttt{complexity\_bias}) & 0.15 \\
\(\beta_0\) (\texttt{cost\_penalty\_beta}) & 0.63 \\
\(\lambda_0\) (\texttt{over\_penalty\_lambda}) & 0.35 \\
\(\alpha\) (\texttt{preference\_power}) & 1.56 \\
\(A_{\mu}^{+}\) (\texttt{max\_mu\_multiplier}) & 3.24 \\
\(A_b^{+}\) (\texttt{max\_bias\_shift}) & 1.05 \\
\(A_{\beta}^{+}\) (\texttt{max\_cost\_relief}) & 13.4 \\
\(A_{\lambda}^{+}\) (\texttt{max\_over\_relief}) & 896 \\
\(A_{\mu}^{-}\) (\texttt{min\_mu\_multiplier}) & 0.377 \\
\(A_b^{-}\) (\texttt{min\_bias\_shift}) & -5.87 \\
\(A_{\beta}^{-}\) (\texttt{min\_cost\_boost}) & 13 \\
\(A_{\lambda}^{-}\) (\texttt{min\_over\_boost}) & 11.9 \\
\bottomrule
\end{tabular}
\end{table}

\subsection{Algorithm Summary}
\label{sec:brick-algorithm-summary}

Given a query \(x\), the deployed router executes the following deterministic
procedure:
\begin{enumerate}
  \item Normalize and truncate \(x\) for the classification models.
  \item Compute capability probabilities \(p(x)\).
  \item Compute complexity label, confidence, and blended \(\tau_q\).
  \item Apply the user preference knob to obtain
  \((\mu(r),b(r),\beta(r),\lambda(r))\).
  \item For every candidate model, compute \(D_m\), \(J_m\), and \(E_m\).
  \item Select the minimum-score model, with deterministic tie-breaking.
\end{enumerate}

Keyword rules may optionally override or bias the capability distribution
before the final score is computed. They do not change the mathematical form of
the scoring objective.

\subsection{Evaluation Metrics}
\label{sec:brick-metrics}

Let \(m_i^\star\) be the router-selected model for query \(i\), and let
\[
  z_{i,m}\in\{0,1\}
\]
indicate whether model \(m\) solved query \(i\). The primary quality metric for
a selected model is
\[
  \mathrm{SelectedAnswerAccuracy}
  =
  \frac{1}{N}\sum_{i=1}^{N} z_{i,m_i^\star}.
\]
This metric answers the deployment question: did the model selected by the
router produce a correct answer?

For economy-focused evaluation, it is also useful to define a cheapest-correct
label
\[
  g_i=\arg\min_{m:z_{i,m}=1} a_m,
\]
with a fallback to the strongest model when no model solves the query. The
exact routing accuracy is
\[
  \mathrm{RouteExactAccuracy}
  =
  \frac{1}{N}\sum_{i=1}^{N} \mathbf{1}[m_i^\star=g_i].
\]
This metric intentionally penalizes over-routing. Therefore it is not the right
objective for the max-performance profile: a stronger but more expensive model
can be correct while still failing the cheapest-correct exact match.

The average routing cost is
\[
  \mathrm{AvgCost}
  =
  \frac{1}{N}\sum_{i=1}^{N} a_{m_i^\star},
\]
and the cost per correct answer is
\[
  \mathrm{CostPerCorrect}
  =
  \frac{\mathrm{AvgCost}}
       {\max(\mathrm{SelectedAnswerAccuracy},\eta)},
\]
where \(\eta\) is a small numerical constant.

\subsection{Locked Full-Dataset Evaluation}
\label{sec:brick-full-dataset-eval}

Table~\ref{tab:brick-full-oof} reports the locked-configuration full-dataset
evaluation over all \(5{,}504\) Dataset~A rows. The protocol is stratified
5-fold out-of-fold: each example is evaluated with skill vectors calibrated on
the other folds, while the knob hyperparameters in
Table~\ref{tab:brick-knob-defaults} remain fixed from prior development. This
uses the whole dataset without per-example skill-vector leakage. It is still an
in-domain calibrated evaluation, not an external blind benchmark.

\begin{table}[ht]
\centering
\caption{Locked aggressive-knob evaluation on all \(5{,}504\) Dataset~A rows
using stratified 5-fold out-of-fold skill calibration. The distribution column
reports selected-model percentages as qwen/deepseek/kimi.}
\label{tab:brick-full-oof}
\small
\begin{tabular}{l r r r l}
\toprule
\textbf{Profile} & \textbf{Selected answer} & \textbf{Route exact} &
\textbf{Avg. cost} & \textbf{Distribution} \\
\midrule
min (\(r=-1\))      & 63.17\% & 63.17\% & 0.100 & 100/0/0 \\
neutral (\(r=0\))   & 68.46\% & 60.10\% & 0.192 & 71/25/3 \\
max (\(r=+1\))      & \textbf{76.98\%} & 17.84\% & 0.537 & 0/31/69 \\
\bottomrule
\end{tabular}
\end{table}

For context, the standalone full-dataset baselines are:
\[
\begin{array}{l r}
\texttt{qwen3.5-9b} & 63.17\%\\
\texttt{deepseek-v4-flash} & 73.69\%\\
\texttt{kimi2.6} & 75.02\%\\
\text{three-model oracle} & 83.25\%.
\end{array}
\]
Thus the max profile exceeds the best single-model baseline in this
in-domain evaluation, while remaining below the oracle ceiling. The gap to the
oracle corresponds to routing opportunities the current feature representation
and decision rule do not yet recover.

\subsection{Scientific Scope}
\label{sec:brick-scientific-scope}

The Brick router is a calibrated decision rule, not a fine-tuned language
model. The calibration step estimates \(S\) and selects routing
hyperparameters; it does not update the weights of
\texttt{qwen3.5-9b}, \texttt{deepseek-v4-flash}, or \texttt{kimi2.6}.

The strongest claim supported by the reported full-dataset experiment is an
in-domain calibrated routing result. The out-of-fold protocol prevents
per-example leakage in skill-vector estimation, but the knob hyperparameters
were developed using Dataset~A. Consequently, this result is appropriate as a
calibrated in-domain routing evaluation and algorithmic ablation. It should not
be described as zero-shot routing, model fine-tuning, or validation on an
independent external benchmark.
